\documentclass[UTF8, 10pt]{ctexart}

\usepackage[margin=0.72in]{geometry}
\usepackage{booktabs}
\usepackage{array}
\usepackage{hyperref}
\usepackage{longtable}
\usepackage{enumitem}
\setlist{leftmargin=*,nosep}
\setlength{\parskip}{0.18em}
\setlength{\parindent}{0pt}
\setlength{\tabcolsep}{4pt}
\renewcommand{\arraystretch}{1.08}

\title{\textbf{CS240：算法设计与分析}\\
课程项目通知}
\author{2026 年春季学期}
\date{主题：现代机器学习与数据应用中的算法}

\begin{document}

\maketitle

\vspace{-1.4em}

\section{时间安排}

\begin{center}
\begin{tabular}{ll}
\toprule
里程碑 & 截止时间 \\
\midrule
项目提案提交 & 5 月 17 日截止 \\
项目展示 & 第 16 周，具体日期待定 \\
最终报告提交 & 项目展示后截止 \\
\bottomrule
\end{tabular}
\end{center}

\section{组队信息与选题}

\subsection{组队人数}

学生可以自由组队，每组最多 \textbf{5 名成员}。

\subsection{选题要求}

每个小组必须从附录中的参考选题中选择一个项目主题，或在获得批准后选择自拟
主题。

\subsection{提交政策}

\begin{itemize}
    \item 项目提案必须通过 Gradescope 提交：
    {\footnotesize\url{https://www.gradescope.com/courses/1270636/assignments/8124428/submissions}}。
    \item 提交时，必须在平台上选择本组队友。
    \item 每个小组只需提交一份提案。
\end{itemize}

\subsection{项目范围}

每个项目都必须展示经典算法如何解决现代问题。项目重点在于理解这种联系，而
不是追求原创研究或最先进的性能表现。

\section{项目阶段与核心要求}

为了保证项目质量，各小组需要完成以下五个阶段。

\subsection{阶段 1：背景与相关工作}

\begin{itemize}
    \item \textbf{背景介绍：} 简要介绍所选择的应用领域。
    \item \textbf{算法总结：} 总结相关的算法形式化建模。
    \item \textbf{文献综述：} 讨论若干代表性论文。
    \item \textbf{选题理由：} 解释为什么所选问题在算法层面上具有研究价值。
    \item \textbf{篇幅要求：} 约 1--2 页，用于展示对问题的深入理解。
\end{itemize}

\subsection{阶段 2：代表性论文 / 方法选择}

\begin{itemize}
    \item \textbf{问题形式化：} 解释所选论文中的正式问题定义。
    \item \textbf{算法描述：} 详细描述核心算法。
    \item \textbf{复杂度分析：} 对算法复杂度进行严谨分析。
    \item \textbf{实验目标：} 明确定义计划复现的实验或结果。
\end{itemize}

\subsection{阶段 3：复现与实现}

\begin{itemize}
    \item \textbf{算法实现：} 正确实现所选算法。
    \item \textbf{结果复现：} 至少复现原论文中部分报告结果。
    \item \textbf{结果分析：} 解释复现结果中的偏差或失败原因。
    \item \textbf{性能分析：} 初步分析运行时间和可扩展性。
\end{itemize}

\subsection{阶段 4：实验评估}

\begin{itemize}
    \item \textbf{运行时间分析：} 测量并分析程序执行时间。
    \item \textbf{可扩展性评估：} 评估算法在数据规模增大时的表现。
    \item \textbf{对比实验：} 将你的实现与至少一个基线方法或变体进行比较。
\end{itemize}

\subsection{阶段 5：可选拓展（加分级别）}

优秀项目可以包含以下拓展内容：

\begin{itemize}
    \item 算法优化或启发式改进。
    \item 近似方法或更好的数据预处理。
    \item 在其他数据集上测试，或开展额外实验。
\end{itemize}

\section{提交内容}

\subsection{A. 项目提案}

项目提案是一份简短文档，约 \textbf{1--2 页}，必须包含以下部分：

\begin{itemize}
    \item \textbf{选题说明：} 所选项目主题、简要动机和应用背景。
    \item \textbf{问题形式化：} 清晰定义计算问题，包括输入 / 输出描述和优化
    目标。
    \item \textbf{相关工作：} 简要总结至少两篇相关论文，并说明具体选择哪篇
    论文或方法进行复现。
    \item \textbf{算法与技术计划：} 涉及的核心算法，例如动态规划或图算法，
    计划采用的实现方式，以及预期使用的数据集或评估设置。
    \item \textbf{项目范围与预期目标：} 小组计划复现的内容，以及可能的可选
    拓展或改进。
    \item \textbf{团队信息：} 小组成员名单和预期分工。
\end{itemize}

\subsection{B. 实验报告}

实验报告应围绕以下项目阶段展开：

\begin{itemize}
    \item \textbf{背景与相关工作：} 介绍应用领域，总结算法形式化建模，并讨论
    代表性论文。
    \item \textbf{代表性论文 / 方法分析：} 解释所选论文的问题形式化、算法描述
    和复杂度分析。
    \item \textbf{复现与实现：} 说明算法实现细节、复现的报告结果，以及任何
    偏差或失败原因。
    \item \textbf{实验评估：} 深入分析运行时间和可扩展性，并与至少一个基线
    方法或变体进行比较。
    \item \textbf{可选拓展（加分级别）：} 记录进一步的优化、启发式改进，或在
    其他数据集上的实验。
\end{itemize}

\subsection{C. 源代码}

源代码提交必须包含：

\begin{itemize}
    \item 核心算法的完整实现。
    \item 如何运行复现实验的文档说明。
\end{itemize}

\section{加分规则}

表现最优秀的小组将获得额外加分。

\begin{center}
\begin{tabular}{lr}
\toprule
排名 & 加分 \\
\midrule
前 3 名小组 & +6 分 \\
前 5 名小组 & +3 分 \\
前 10 名小组 & +1 分 \\
\bottomrule
\end{tabular}
\end{center}

\section*{附录：参考选题}
\addcontentsline{toc}{section}{附录：参考选题}

\small
\begin{longtable}{>{\raggedleft\arraybackslash}p{0.05\linewidth}
                  p{0.20\linewidth}
                  p{0.36\linewidth}
                  p{0.29\linewidth}}
\toprule
\# & 主题 & 相关概念 & 应用场景 \\
\midrule
\endfirsthead
\toprule
\# & 主题 & 相关概念 & 应用场景 \\
\midrule
\endhead
1 & 知识图谱摘要 & 顶点覆盖、集合覆盖、近似算法、贪心方法
  & 知识图谱、搜索引擎、信息检索 \\
2 & 社区发现 / 图划分 & 最小割、最大流、图划分、谱方法
  & 社交网络、推荐系统 \\
3 & 影响力最大化 & 贪心算法、近似算法、子模优化
  & 病毒式营销、虚假信息分析 \\
4 & 恶意软件传播控制 & 支配集、顶点覆盖、最短路径、网络优化
  & 网络安全、网络防御 \\
5 & 配送路线优化 & 旅行商问题、最小生成树、近似算法、贪心启发式
  & 物流、外卖配送、交通运输 \\
6 & 二分图匹配 & 二分图匹配、匈牙利算法、网络流
  & 推荐系统、工作分配 \\
7 & 虚假评论 / 欺诈检测 & 团检测、稠密子图、图挖掘、NP 完全问题
  & 欺诈检测、评论垃圾信息分析 \\
8 & 翻译中的编辑距离 & 动态规划、序列比对、字符串算法
  & 机器翻译、拼写检查、NLP 评估 \\
9 & DNA / 蛋白质序列比对 & 最长公共子序列、动态规划、局部 / 全局比对
  & 生物信息学、计算生物学 \\
10 & 字幕或转录文本对齐 & 动态规划、近似字符串匹配、序列比对
   & 语音系统、视频处理 \\
11 & LLM 上下文选择 & 背包问题、集合覆盖、近似算法
   & 检索增强生成（RAG）、LLM 系统 \\
12 & 对抗性路径规划 & 最短路径、A* 搜索、动态规划、启发式搜索
   & 对抗机器学习、安全系统 \\
13 & 内容感知图像缩放 & 动态规划、最短路径、贪心优化
   & 在保留重要视觉内容的同时调整图像大小 \\
14 & 基于图割的图像分割 & 最小割 / 最大流、图构建、能量最小化
   & 图像中的目标分割 \\
\bottomrule
\end{longtable}
\normalsize

\end{document}

